%%%%%%%%%%%%%%%%%%%%%%%%%%%%%%%%%%%%%%%%%%%%%%%%%%%%%%%%%%%%%%%%%%%%%%%%%%%%%%%
% CHAPTER 12: WILLINGNESS TO ACT
%%%%%%%%%%%%%%%%%%%%%%%%%%%%%%%%%%%%%%%%%%%%%%%%%%%%%%%%%%%%%%%%%%%%%%%%%%%%%%%
% Version: 2.0 (January 2026)
%
% Purpose: Introduces the Willingness function WAX and threshold θ that
%          determine whether effective utility translates into action.
%          Answers the seventh CORE question: "How ready am I to act?"
%
% Primary Appendix: AV - CORE-READY (The Willingness Function)
% Secondary Appendices: AU (CORE-AWARE), C (CORE-WHAT), V (CORE-WHEN), AW (CORE-STAGE)
%
% Prerequisites: Ch. 10 (Welfare/FEPSDE), Ch. 11 (Awareness)
% Leads to: Ch. 13 (Behavioral Change Journey)
%
% Chapter Type: A (CORE Chapter - answers READY question)
%
% EBF Questions addressed:
%   Question 1: HOW does he/she think? (φ)
%   Question 4: HOW LARGE is WAX and WHAT THRESHOLDS exist? (WAX, θ)
%%%%%%%%%%%%%%%%%%%%%%%%%%%%%%%%%%%%%%%%%%%%%%%%%%%%%%%%%%%%%%%%%%%%%%%%%%%%%%%

\chapter{Willingness to Act: From Thinking to Deciding}
\label{ch:willingness}

%==============================================================================
% DIE 9 FRAGEN NAVIGATION BOX
%==============================================================================
\BCMQuestionNav{12}

%==============================================================================
% 9C CORE CONNECTION BOX
%==============================================================================
\begin{tcolorbox}[
    colback=violet!5!white,
    colframe=violet!75!black,
    title={\textbf{9C CORE Connection: READY (Question 7)}},
    fonttitle=\bfseries
]
\textbf{The 7th CORE Question:} How ready am I to act?

\medskip
\begin{tabular}{ll}
\textbf{CORE Appendix:} & AV (CORE-READY) \\
\textbf{Output:} & Willingness $WAX$, Threshold $\theta$, Gap $\Delta = WAX - \theta$ \\
\textbf{Key Insight:} & Awareness is necessary but not sufficient; action requires $WAX > \theta$ \\
\textbf{Novel Contribution:} & 99 axioms, $\varphi$-decomposition, threshold dynamics
\end{tabular}

\medskip
\textbf{The Action Condition:}
\begin{equation}
\boxed{\text{Action} \Leftrightarrow WAX \geq \theta}
\label{eq:action-condition}
\end{equation}

\textit{AWARE (AU) transforms potential into effective utility. READY (AV) determines whether effective utility triggers action. Together they explain the intention-behavior gap.}
\end{tcolorbox}

%==============================================================================
% APPENDIX REFERENCES BOX
%==============================================================================
\begin{tcolorbox}[
    colback=yellow!10!white,
    colframe=orange!75!black,
    title={\textbf{Appendix References for This Chapter}},
    fonttitle=\bfseries
]
\begin{tabular}{lll}
\textbf{Appendix} & \textbf{Content} & \textbf{Section} \\
\midrule
AV (CORE-READY) & 99 axioms (WAX-A1 to WAX-A103) & All sections \\
AV Section 2 & Phi Decomposition ($\varphi$) & \ref{sec:12-phi} \\
AV Section 3 & WAX Master Equation & \ref{sec:12-wax} \\
AV Section 4 & Threshold Dynamics ($\theta$) & \ref{sec:12-theta} \\
AU (CORE-AWARE) & Effective utility $U^* = A(\cdot) \times U_{pot}$ & \ref{sec:12-intro} \\
G (REF-GLOSSARY) & Symbol definitions & All sections
\end{tabular}
\end{tcolorbox}


%%%%%%%%%%%%%%%%%%%%%%%%%%%%%%%%%%%%%%%%%%%%%%%%%%%%%%%%%%%%%%%%%%%%%%%%%%%%%%%
%                          PART I: INTRODUCTION
%%%%%%%%%%%%%%%%%%%%%%%%%%%%%%%%%%%%%%%%%%%%%%%%%%%%%%%%%%%%%%%%%%%%%%%%%%%%%%%

\section{Introduction: The Action Condition}
\label{sec:12-intro}

%------------------------------------------------------------------------------

Chapter 11 established \textit{Awareness} (AWX)---what people perceive at the
Moment of Truth. But perception alone does not guarantee action. This chapter
addresses the next critical question: \textbf{Willingness to Act} (WAX).

\begin{intuition}[The Knowing-Doing Gap]
Consider Thomas, who knows exactly that:
\begin{itemize}[nosep]
  \item Exercise would improve his health ($U^P$ high)
  \item He would feel better emotionally ($U^E$ high)
  \item His doctor recommends it ($AWX$ high---he's aware)
\end{itemize}
Yet Thomas hasn't exercised in months. Why?

The gap between \textbf{knowing} and \textbf{doing} is where Willingness operates.
Thomas is aware but not yet \textbf{willing}---his $WAX < \theta$.
\end{intuition}

%------------------------------------------------------------------------------
% CENTRAL QUESTION BOX
%------------------------------------------------------------------------------
\paragraph{The Central Question.}
This chapter addresses the 7th of the 9 BCM Questions:

\begin{center}
\fcolorbox{violet!75!black}{violet!5!white}{\parbox{0.85\textwidth}{
\textbf{Question 7 (READY):} How ready am I to act on what I'm aware of?
}}
\end{center}

\begin{definition}[Willingness to Act]
\textbf{Willingness to Act (WAX)} is the degree to which an individual is prepared
to execute a behavior, given their awareness of its consequences. WAX integrates:
\begin{itemize}[nosep]
\item How the person \textit{thinks} about the decision (thinking style $\varphi$)
\item The perceived utilities across all dimensions ($U^* = AWX \cdot U$)
\item The barriers and facilitators in the decision context
\end{itemize}
\end{definition}

\begin{intuition}[Awareness vs. Willingness]
\begin{itemize}[nosep]
\item \textbf{Awareness} answers: ``Do I \textit{see} the benefits?''
\item \textbf{Willingness} answers: ``Am I \textit{ready} to act on them?''
\end{itemize}
A person can be fully aware of the benefits of exercise but still unwilling to
go to the gym. The gap between knowing and doing is where Willingness operates.
\end{intuition}

%------------------------------------------------------------------------------
% CHAPTER OVERVIEW
%------------------------------------------------------------------------------
\paragraph{Chapter Overview.}
This chapter proceeds as follows:
\begin{itemize}
    \item Section \ref{sec:12-phi}: The thinking style parameter $\varphi$
    \item Section \ref{sec:12-wax}: The WAX formula and calculation
    \item Section \ref{sec:12-theta}: Thresholds and barrier dynamics
    \item Section \ref{sec:12-comparison}: The critical comparison $WAX$ vs. $\theta$
    \item Section \ref{sec:12-anna}: Complete worked example (Anna)
    \item Section \ref{sec:12-integration}: Integration with 9C framework
    \item Section \ref{sec:12-summary}: Summary and transition
\end{itemize}


%%%%%%%%%%%%%%%%%%%%%%%%%%%%%%%%%%%%%%%%%%%%%%%%%%%%%%%%%%%%%%%%%%%%%%%%%%%%%%%
%                     PART II: THE THINKING STYLE
%%%%%%%%%%%%%%%%%%%%%%%%%%%%%%%%%%%%%%%%%%%%%%%%%%%%%%%%%%%%%%%%%%%%%%%%%%%%%%%

\section{Question 1: How Does the Person Think? ($\varphi$)}
\label{sec:12-phi}

%------------------------------------------------------------------------------

The first question this chapter answers is fundamental:
\textbf{How does the person think about multi-dimensional decisions?}

\subsection{The Thinking Style Parameter $\varphi$}

\begin{definition}[Thinking Style]
\begin{equation}
\varphi \in [0, 1]
\label{eq:phi-range}
\end{equation}
The thinking style $\varphi$ determines how a person integrates multiple
utility dimensions into an overall evaluation.
\end{definition}

\begin{intuition}[What is $\varphi$?]
When facing a decision with multiple dimensions (financial, social, environmental,
identity...), people think in fundamentally different ways:

\begin{center}
\begin{tabular}{cp{10cm}}
\toprule
$\varphi \to 1$ & \textbf{``Connected Thinking'' (Complementary)}\\
& ``I only act if EVERYTHING is right''\\
& The weakest dimension dominates the decision\\
& Like a chain: only as strong as the weakest link\\
& Example: ``I won't buy unless price AND environment AND social proof all check out''\\[8pt]

$\varphi \to 0$ & \textbf{``Separated Thinking'' (Additive)}\\
& ``I look at the total benefit''\\
& All dimensions add up; strengths compensate for weaknesses\\
& Example: ``Great financial savings can make up for weak environmental benefits''\\
\bottomrule
\end{tabular}
\end{center}
\end{intuition}

\subsection{$\varphi$-Decomposition}

The thinking style is not fixed---it depends on both the person and the context:

\begin{definition}[Phi Decomposition]
\begin{equation}
\varphi = \varphi_0 + \Delta\varphi(\Psi)
\label{eq:phi-decomposition}
\end{equation}
where:
\begin{itemize}[nosep]
\item $\varphi_0$ = baseline thinking style (personality trait, relatively stable)
\item $\Delta\varphi(\Psi) = \sum_j \gamma_j \cdot \Psi_j$ = context-induced shift
\end{itemize}
\end{definition}

\begin{intuition}[What shifts $\varphi$?]
\begin{center}
\begin{tabular}{lcp{6cm}}
\toprule
\textbf{Context} & \textbf{Effect on $\varphi$} & \textbf{Why}\\
\midrule
High uncertainty ($\Psi_{UNS}$) & $\downarrow$ & Fall back to self-interest\\
Strong group identity ($\Psi_{SEG}$) & $\uparrow$ & Norms become binding\\
Salient identity ($\Psi_{IDN}$) & $\uparrow$ & ``Who am I?'' dominates\\
Social observation ($\Psi_{SOC}$) & $\uparrow$ & Others are watching\\
\bottomrule
\end{tabular}
\end{center}

\textbf{Example:} Anna normally thinks additively ($\varphi_0 = 0.35$). But when
her neighbor mentions ``everyone on our street is getting heat pumps,'' her
$\varphi$ shifts upward (+0.20) because social norms become salient.
\end{intuition}

\begin{cross-reference}
\textbf{Formal Specification:} See Appendix AV (CORE-READY), Section 2 ``Phi Decomposition
and Estimation'' for the complete $\gamma$-matrix and LLM-MC estimation methodology.
\end{cross-reference}


%%%%%%%%%%%%%%%%%%%%%%%%%%%%%%%%%%%%%%%%%%%%%%%%%%%%%%%%%%%%%%%%%%%%%%%%%%%%%%%
%                     PART III: THE WAX FORMULA
%%%%%%%%%%%%%%%%%%%%%%%%%%%%%%%%%%%%%%%%%%%%%%%%%%%%%%%%%%%%%%%%%%%%%%%%%%%%%%%

\section{Question 4a: How Large is the Willingness? (WAX)}
\label{sec:12-wax}

%------------------------------------------------------------------------------

\subsection{The WAX Master Equation}

\begin{definition}[Willingness to Act]
\begin{equation}
WAX = \varphi \cdot \min_i\{U_i^*\} + (1-\varphi) \cdot \sum_i U_i^*
\label{eq:wax-master}
\end{equation}
where $U_i^* = AWX_i \cdot U_i$ is the \textit{effective utility}---the utility
the person actually perceives.
\end{definition}

\begin{intuition}[What does this formula mean?]
The formula is a weighted combination of two ways of evaluating options:

\begin{enumerate}
\item \textbf{Complementary logic} ($\min$): The weakest dimension determines value\\
$\rightarrow$ ``A meal is only as good as its worst dish''

\item \textbf{Additive logic} ($\sum$): All dimensions contribute to total value\\
$\rightarrow$ ``A great dessert can compensate for mediocre soup''
\end{enumerate}

$\varphi$ determines the weight between these two logics.
\end{intuition}

\subsection{Worked Example: Anna's WAX Calculation}

\begin{example}[Anna's Willingness]
Anna is considering a heat pump. Her values:

\begin{center}
\begin{tabular}{lcccc}
\toprule
\textbf{Dimension} & $U$ (objective) & $AWX$ & $U^* = AWX \cdot U$ \\
\midrule
Individual (INU) & 0.7 & 0.5 & 0.35 \\
Collective (KNU) & 0.6 & 0.5 & 0.30 \\
Identity (IDN) & 0.5 & 0.5 & 0.25 \\
\bottomrule
\end{tabular}
\end{center}

With $\varphi = 0.35$:
\begin{align}
\min\{U^*\} &= \min\{0.35, 0.30, 0.25\} = 0.25\\
\sum U^* &= 0.35 + 0.30 + 0.25 = 0.90\\
WAX &= 0.35 \cdot 0.25 + 0.65 \cdot 0.90\\
&= 0.088 + 0.585 = \mathbf{0.67}
\end{align}
\end{example}

\begin{cross-reference}
\textbf{Formal Specification:} See Appendix AV (CORE-READY), Section 3 ``The WAX Master
Equation'' for the complete static and dynamic formulation.
\end{cross-reference}


%%%%%%%%%%%%%%%%%%%%%%%%%%%%%%%%%%%%%%%%%%%%%%%%%%%%%%%%%%%%%%%%%%%%%%%%%%%%%%%
%                     PART IV: THRESHOLDS
%%%%%%%%%%%%%%%%%%%%%%%%%%%%%%%%%%%%%%%%%%%%%%%%%%%%%%%%%%%%%%%%%%%%%%%%%%%%%%%

\section{Question 4b: What Are the Thresholds? ($\theta$)}
\label{sec:12-theta}

%------------------------------------------------------------------------------

Willingness alone is not sufficient. It must \textit{exceed a threshold} for
action to occur.

\subsection{The Threshold Formula}

\begin{definition}[Threshold Decomposition]
\begin{equation}
\theta = \theta_{base}(d) + \sum_j \delta_j \cdot \Psi_j
\label{eq:theta-decomposition}
\end{equation}
where $\theta_{base}(d)$ depends on decision type and $\delta_j$ are context
sensitivity coefficients.
\end{definition}

\begin{intuition}[What determines the threshold?]

\textbf{1. Base difficulty} depending on decision type:

\begin{center}
\begin{tabular}{lcp{5cm}}
\toprule
\textbf{Decision Type} & $\theta_{base}$ & \textbf{Example}\\
\midrule
Continue habit & 0.2 & Keep using gas heating\\
Small adjustment & 0.5 & Switch electricity provider\\
Moderate change & 0.8 & Buy an e-bike\\
Major investment & 1.2 & Install heat pump\\
Life-changing & 1.5+ & Move abroad for climate\\
\bottomrule
\end{tabular}
\end{center}

\textbf{2. Context adjustments} ($\sum \delta_j \cdot \Psi_j$):
\begin{itemize}[nosep]
\item High uncertainty $\rightarrow$ threshold goes UP (harder to decide)
\item Good subsidies ($\Psi_{REG}$) $\rightarrow$ threshold goes DOWN (easier)
\item Social proof ($\Psi_{SOC}$) $\rightarrow$ threshold goes DOWN
\end{itemize}
\end{intuition}

\begin{cross-reference}
\textbf{Formal Specification:} See Appendix AV (CORE-READY), Axiom WAX-A103
``Threshold Decomposition'' for the complete $\delta$-matrix and decision type classification.
\end{cross-reference}


%%%%%%%%%%%%%%%%%%%%%%%%%%%%%%%%%%%%%%%%%%%%%%%%%%%%%%%%%%%%%%%%%%%%%%%%%%%%%%%
%                     PART V: THE CRITICAL COMPARISON
%%%%%%%%%%%%%%%%%%%%%%%%%%%%%%%%%%%%%%%%%%%%%%%%%%%%%%%%%%%%%%%%%%%%%%%%%%%%%%%

\section{The Critical Comparison: WAX vs. $\theta$}
\label{sec:12-comparison}

%------------------------------------------------------------------------------

The entire EBF framework culminates in one comparison:

\begin{center}
\fcolorbox{violet!75!black}{violet!5!white}{\parbox{0.5\textwidth}{
\centering
\Large
$WAX \gtrless \theta$
\\[6pt]
\normalsize
\textit{Does willingness exceed the threshold?}
}}
\end{center}

\subsection{The Three Regimes}

\begin{table}[h]
\centering
\caption{The Three Decision Regimes}
\label{tab:three-regimes}
\begin{tabular}{cccp{5cm}}
\toprule
\textbf{Regime} & \textbf{Condition} & \textbf{P(Action)} & \textbf{Plain Language}\\
\midrule
\cellcolor{red!15} Below & $WAX < \theta$ & $< 50\%$ & ``I don't want it enough''\\
\cellcolor{yellow!15} Critical & $WAX \approx \theta$ & $\approx 50\%$ & ``I'm on the fence''\\
\cellcolor{green!15} Above & $WAX > \theta$ & $> 50\%$ & ``Yes, I'll do it''\\
\bottomrule
\end{tabular}
\end{table}

\subsection{The Gap: $\Delta = WAX - \theta$}

\begin{definition}[The Willingness Gap]
\begin{equation}
\Delta = WAX - \theta
\label{eq:gap}
\end{equation}
The gap measures how far from action threshold an individual stands.
\end{definition}

\begin{intuition}[Interpreting the Gap]
\begin{center}
\begin{tabular}{rcp{6cm}}
$\Delta > 0$ & $\checkmark$ & Ready to act\\
$\Delta < 0$ & $\times$ & Not ready\\
$\Delta \approx 0$ & $\sim$ & Tipping point (Moment of Truth)\\
\end{tabular}
\end{center}

\textbf{Intervention implication:}
\begin{itemize}[nosep]
\item Small gap ($|\Delta| < 0.2$): Light nudge sufficient
\item Medium gap ($|\Delta| = 0.2$--$0.5$): Substantive intervention needed
\item Large gap ($|\Delta| > 0.5$): Structural change required
\end{itemize}
\end{intuition}

\subsection{Two Paths to Close the Gap}

To move someone from ``no'' to ``yes,'' there are two complementary paths:

\begin{table}[h]
\centering
\caption{Two Paths to Behavioral Change}
\label{tab:two-paths}
\begin{tabular}{lp{5cm}p{5cm}}
\toprule
& \textbf{Path 1: Increase WAX} & \textbf{Path 2: Decrease $\theta$}\\
& ``Make it more attractive'' & ``Make it easier''\\
\midrule
Levers & $\uparrow \varphi$ (shift mindset) & $\downarrow \theta_{base}$ (simplify decision)\\
& $\uparrow AWX$ (increase awareness) & $\downarrow$ friction (remove barriers)\\
& $\uparrow U$ (enhance benefits) & Change defaults (opt-out vs. opt-in)\\
\midrule
Example & ``Your neighbors got one too!'' & ``One-stop-shop installation''\\
& ``You'll save CHF 2,000/year'' & ``Pre-filled application form''\\
\bottomrule
\end{tabular}
\end{table}

\textbf{Best practice:} Combine both paths for maximum effect.

\begin{cross-reference}
\textbf{Formal Specification:} See Appendix AV (CORE-READY), Section 5 ``The Critical
Comparison: WAX vs. $\theta$'' for complete intervention strategy framework.
\end{cross-reference}


%%%%%%%%%%%%%%%%%%%%%%%%%%%%%%%%%%%%%%%%%%%%%%%%%%%%%%%%%%%%%%%%%%%%%%%%%%%%%%%
%                     PART VI: COMPLETE EXAMPLE
%%%%%%%%%%%%%%%%%%%%%%%%%%%%%%%%%%%%%%%%%%%%%%%%%%%%%%%%%%%%%%%%%%%%%%%%%%%%%%%

\section{Complete Example: Anna's Journey}
\label{sec:12-anna}

%------------------------------------------------------------------------------

\begin{example}[Anna: From 34\% to 54\%]

\textbf{BASELINE (before intervention):}
\begin{align}
\varphi &= 0.35 \quad \text{(mostly additive)}\\
AWX &= 0.5 \quad \text{(sees only half the benefits)}\\
WAX &= 0.67\\
\theta &= 1.32 \quad \text{(major investment + uncertainty)}\\
\Delta &= 0.67 - 1.32 = \mathbf{-0.65}\\
P &= \sigma(-0.65) \approx \mathbf{34\%}
\end{align}

\textbf{INTERVENTION:}
\begin{enumerate}
\item Show concrete savings calculation $\rightarrow$ $AWX$: $0.5 \to 0.8$
\item ``The Müllers got one!'' (social norm) $\rightarrow$ $\varphi$: $0.35 \to 0.55$
\item Highlight subsidies $\rightarrow$ $\theta$: $-0.30$
\item Simplified application $\rightarrow$ $\theta$: $-0.26$
\end{enumerate}

\textbf{AFTER INTERVENTION:}
\begin{align}
\varphi &= 0.55\\
AWX &= 0.8\\
WAX &= 0.90\\
\theta &= 0.76\\
\Delta &= 0.90 - 0.76 = \mathbf{+0.14}\\
P &= \sigma(+0.14) \approx \mathbf{54\%}
\end{align}

\textbf{Result:} Anna crossed the tipping point. Probability increased by +20 percentage points.
\end{example}

\begin{intuition}[What Made the Difference?]
Anna's journey illustrates the power of combined interventions:
\begin{itemize}[nosep]
  \item \textbf{Path 1 (WAX $\uparrow$):} Awareness and social norms increased WAX from 0.67 to 0.90
  \item \textbf{Path 2 ($\theta$ $\downarrow$):} Subsidies and simplification reduced θ from 1.32 to 0.76
  \item \textbf{Combined effect:} Gap changed from $-0.65$ to $+0.14$
\end{itemize}
Neither path alone would have been sufficient.
\end{intuition}


%%%%%%%%%%%%%%%%%%%%%%%%%%%%%%%%%%%%%%%%%%%%%%%%%%%%%%%%%%%%%%%%%%%%%%%%%%%%%%%
%                     PART VII: 9C INTEGRATION
%%%%%%%%%%%%%%%%%%%%%%%%%%%%%%%%%%%%%%%%%%%%%%%%%%%%%%%%%%%%%%%%%%%%%%%%%%%%%%%

\section{Integration with the CORE Framework}
\label{sec:12-integration}

%------------------------------------------------------------------------------

\subsection{Connection to CORE-AWARE (Awareness)}

\begin{cross-reference}
The Awareness function $A(\cdot)$ from CORE-AWARE provides the input to WAX:
\begin{itemize}[nosep]
  \item Effective utility $U^* = A(\cdot) \times U_{pot}$ feeds into WAX formula
  \item AWX and WAX are sequential: first perceive, then be willing
  \item Low awareness ($A \approx 0$) prevents willingness from forming
\end{itemize}
See Chapter 11 and Appendix AU (CORE-AWARE) for the awareness specification.
\end{cross-reference}

\subsection{Connection to CORE-WHAT (FEPSDE)}

\begin{cross-reference}
The FEPSDE dimensions from CORE-WHAT define the utility components:
\begin{itemize}[nosep]
  \item WAX operates on all 144 utility components
  \item $\min\{U^*\}$ and $\sum U^*$ computed across INU, KNU, IDN
  \item Different dimensions have different threshold effects
\end{itemize}
See Chapter 10 and Appendix C (CORE-WHAT) for the utility specification.
\end{cross-reference}

\subsection{Connection to CORE-WHEN (Context)}

\begin{cross-reference}
The context vector $\Psi$ from CORE-WHEN modulates both $\varphi$ and $\theta$:
\begin{itemize}[nosep]
  \item $\varphi = \varphi_0 + \sum_j \gamma_j \cdot \Psi_j$ (thinking style shift)
  \item $\theta = \theta_{base} + \sum_j \delta_j \cdot \Psi_j$ (threshold adjustment)
  \item Context changes can shift both WAX and $\theta$ simultaneously
\end{itemize}
See Chapter 9 and Appendix V (CORE-WHEN) for the context specification.
\end{cross-reference}

\subsection{Connection to CORE-STAGE (BCJ)}

\begin{cross-reference}
The Behavioral Change Journey from CORE-STAGE interacts with WAX and $\theta$:
\begin{itemize}[nosep]
  \item Stage $S_4$ (Intending): $WAX \approx \theta$---on the fence
  \item Stage $S_5$ (Acting): $WAX > \theta$---threshold crossed
  \item Stage $S_6$ (Maintaining): $\theta \to 0$---automatic behavior
\end{itemize}
See Chapter 13 and Appendix AW (CORE-STAGE) for the BCJ specification.
\end{cross-reference}

%------------------------------------------------------------------------------
\subsection{The 9C Integration Table}

\begin{table}[h]
\centering
\begin{tabular}{lll}
\toprule
\textbf{CORE} & \textbf{Question} & \textbf{Willingness Role} \\
\midrule
WHO (AAA) & Who has utility? & WAX varies by level $L$ \\
WHAT (C) & What is utility? & WAX operates on 144 components \\
HOW (B) & How interact? & $\varphi$ captures complementarity \\
WHEN (V) & When context? & Context modulates $\varphi$ and $\theta$ \\
WHERE (BBB) & Where parameters? & 99 axiom parameters \\
AWARE (AU) & How aware? & $U^* = A(\cdot) \times U_{pot}$ feeds WAX \\
\textbf{READY (AV)} & \textbf{How willing?} & \textbf{$WAX > \theta$ triggers action} \\
STAGE (AW) & Where in journey? & Stage affects $\theta$ level \\
\bottomrule
\end{tabular}
\caption{Integration of READY with the 9C framework}
\label{tab:8c-ready-integration}
\end{table}


%%%%%%%%%%%%%%%%%%%%%%%%%%%%%%%%%%%%%%%%%%%%%%%%%%%%%%%%%%%%%%%%%%%%%%%%%%%%%%%
%                     PART VIII: AXIOM OVERVIEW
%%%%%%%%%%%%%%%%%%%%%%%%%%%%%%%%%%%%%%%%%%%%%%%%%%%%%%%%%%%%%%%%%%%%%%%%%%%%%%%

\section{Willingness Axioms: Overview}
\label{sec:12-axioms}

%------------------------------------------------------------------------------

Appendix AV provides the complete formal specification of all 99+ Willingness axioms.

\begin{table}[h]
\centering
\caption{Willingness Axiom Categories}
\label{tab:axiom-categories}
\begin{tabular}{lcp{6cm}}
\toprule
\textbf{Category} & \textbf{Axioms} & \textbf{Key Content}\\
\midrule
Fundamental & WAX-A1 to A18 & Core definitions, $\varphi$, threshold\\
Barrier & WAX-A19 to A40 & What reduces willingness\\
Facilitator & WAX-A41 to A60 & What increases willingness\\
Dynamic & WAX-A61 to A80 & How willingness changes over time\\
Context & WAX-A81 to A99 & $\Psi$-modulation effects\\
Mapping & WAX-A100 to A103 & Market integration, kinship, thresholds\\
\bottomrule
\end{tabular}
\end{table}

\begin{cross-reference}
\textbf{Complete Formal Specification:} See Appendix AV (CORE-READY) for all axioms
with 8-factor documentation (ID, Name, Formal Statement, $\Psi$-Sensitivity,
Empirical Reference, Testable Prediction, Interactions, Intervention Lever).
\end{cross-reference}


%%%%%%%%%%%%%%%%%%%%%%%%%%%%%%%%%%%%%%%%%%%%%%%%%%%%%%%%%%%%%%%%%%%%%%%%%%%%%%%
%                          PART IX: SUMMARY
%%%%%%%%%%%%%%%%%%%%%%%%%%%%%%%%%%%%%%%%%%%%%%%%%%%%%%%%%%%%%%%%%%%%%%%%%%%%%%%

\section{Summary}
\label{sec:12-summary}

%------------------------------------------------------------------------------

This chapter has established:

\begin{enumerate}
    \item \textbf{Thinking Style ($\varphi$):} How people integrate multiple utility
          dimensions---from additive ($\varphi = 0$) to complementary ($\varphi = 1$),
          with context-dependent shifts.

    \item \textbf{WAX Formula:} $WAX = \varphi \cdot \min\{U^*\} + (1-\varphi) \cdot \sum U^*$
          ---the weighted combination of chain logic and sum logic.

    \item \textbf{Threshold Dynamics:} $\theta = \theta_{base}(d) + \sum_j \delta_j \cdot \Psi_j$
          ---base difficulty adjusted by context.

    \item \textbf{The Critical Comparison:} Action occurs when $WAX > \theta$;
          the gap $\Delta = WAX - \theta$ determines intervention intensity.

    \item \textbf{Two Paths to Change:} Increase WAX (make attractive) or
          decrease $\theta$ (make easy)---best when combined.

    \item \textbf{99 Axioms:} Complete formal specification in Appendix AV
          covering barriers, facilitators, dynamics, and context effects.
\end{enumerate}

%------------------------------------------------------------------------------
\paragraph{Formal Details.}
The complete formal treatment appears in Appendix~AV (CORE-READY), which contains:
\begin{itemize}[nosep]
    \item 99 axioms in 8-factor documentation format
    \item $\varphi$-decomposition with $\gamma$-matrix
    \item Threshold decomposition with $\delta$-matrix
    \item LLM-MC estimation methodology
    \item 94+ scientific references
\end{itemize}

%------------------------------------------------------------------------------
\paragraph{What Comes Next.}
Chapter 13 introduces the \textbf{Behavioral Change Journey (BCJ)}, which tracks
where an individual stands in the process of change. The BCJ explains why the
same intervention works at one stage but fails at another---and how $\theta$
evolves as behavior becomes habitual.

%------------------------------------------------------------------------------
\begin{cross-reference}
\textbf{Reading Path:}
\begin{itemize}
    \item For complete axiom specification: Appendix AV (CORE-READY)
    \item For $\varphi$ estimation: Appendix AV, Section 2
    \item For threshold dynamics: Appendix AV, Section 4
    \item For awareness mechanics: Chapter 11 + Appendix AU (CORE-AWARE)
    \item For BCJ (next): Chapter 13 + Appendix AW (CORE-STAGE)
\end{itemize}
\end{cross-reference}

%%%%%%%%%%%%%%%%%%%%%%%%%%%%%%%%%%%%%%%%%%%%%%%%%%%%%%%%%%%%%%%%%%%%%%%%%%%%%%%
% END OF CHAPTER 12
%%%%%%%%%%%%%%%%%%%%%%%%%%%%%%%%%%%%%%%%%%%%%%%%%%%%%%%%%%%%%%%%%%%%%%%%%%%%%%%
